\clearpage
\section*{Exercise 3: Create STM32 Project}
\addcontentsline{toc}{section}{Exercise 3: Create STM32 Project}

\begin{questionbox}{1. Problem Statement}
Create an STM32 project whose GPIO assignments match the Proteus design, and configure a timer interrupt with an interval of approximately 10 ms.
\end{questionbox}

\vspace{0.4cm}
\textbf{2. Project Configuration}

The supplied project targets an STM32F1 device and uses the STM32 HAL. The clock configuration selects the internal high-speed oscillator (HSI) at 8 MHz without enabling the PLL. GPIO initialization enables ports A and B, configures the traffic-light and display signals as outputs, and configures PA13--PA15 as active-low button inputs with pull-ups.

The timer setup is defined in \texttt{MX\_TIM2\_Init()}:

\begin{lstlisting}[language=C, caption={TIM2 prescaler and period settings}, label={lst:exercise3_tim2}, captionpos=b]
htim2.Instance = TIM2;
htim2.Init.Prescaler = 7999;
htim2.Init.CounterMode = TIM_COUNTERMODE_UP;
htim2.Init.Period = 9;
htim2.Init.ClockDivision = TIM_CLOCKDIVISION_DIV1;
\end{lstlisting}

\textbf{3. Timer Calculation}

Assuming the TIM2 input clock is 8 MHz, the prescaler and auto-reload values give the following update interval:

\[
 f_{\mathrm{counter}} = \frac{8\,\mathrm{MHz}}{7999+1} = 1\,\mathrm{kHz}
\]
\[
 f_{\mathrm{update}} = \frac{1\,\mathrm{kHz}}{9+1} = 100\,\mathrm{Hz}
 \qquad\Longrightarrow\qquad
 T_{\mathrm{update}} = 10\,\mathrm{ms}.
\]

The main loop starts the timer in interrupt mode with \texttt{HAL\_TIM\_Base\_Start\_IT()}, while \texttt{HAL\_TIM\_PeriodElapsedCallback()} advances the software tick counter and samples the buttons. The timer callback is therefore the shared time base for button filtering, traffic-phase timing, and LED blinking.

\textbf{4. Implementation Notes}

\begin{itemize}[leftmargin=1.2cm]
    \item \texttt{main.h} centralizes the GPIO pin definitions so the implementation uses consistent signal names.
    \item The seven-segment display is multiplexed: PA8--PA11 select one digit at a time, while the corresponding segment bus is written on GPIOB.
    \item The uploaded \texttt{main.c} does not include the HAL MSP initialization or the startup interrupt-handler file. The CubeIDE project must also enable the TIM2 peripheral clock and NVIC interrupt, and route \texttt{TIM2\_IRQHandler()} to \texttt{HAL\_TIM\_IRQHandler(\&htim2)}.
\end{itemize}

% IMAGE SLOT: Add the screenshot described below and save it as ex03_cubeide_tim2_configuration.png in the images/ folder.
% If the file is absent, a labeled placeholder is displayed; the document remains compilable.
\begin{figure}[H]
    \centering
    \IfFileExists{images/ex03_cubeide_tim2_configuration.png}{%
        \includegraphics[width=0.98\textwidth]{images/ex03_cubeide_tim2_configuration.png}%
    }{%
        \fbox{\parbox[c][4.3cm][c]{0.92\textwidth}{%
            \centering\textbf{IMAGE SLOT -- STM32CubeIDE project and TIM2 settings}\\[6pt]
            Insert a screenshot from STM32CubeIDE showing the project/source file and the TIM2 prescaler/period values (PSC = 7999, ARR = 9), or a successful build view alongside those settings.\\[6pt]
            \texttt{images/ex03\_cubeide\_tim2\_configuration.png}%
        }}%
    }
    \caption{STM32CubeIDE setup used for the 10 ms TIM2 update}
    \label{fig:exercise3_cubeide_tim2}
\end{figure}

\begin{center}
    \href{https://github.com/haunguyen-2006/STM32LAB1}{\textbf{GitHub -- STM32LAB1 Source Repository}}
\end{center}

\begin{resultbox}
With an 8 MHz timer input clock, PSC = 7999 and ARR = 9 produce the required 10 ms timer update. This calculation depends on the actual timer clock and on the TIM2 clock/NVIC setup in the rest of the STM32 project.
\end{resultbox}
